\SourceLeaf{35}

The value found for the fraction $\dfrac{P_\theta}{P_{\theta+1}}$ is therefore too small
if we take the first two terms in the numerator and only the first in the
denominator. In other words, the true value of that fraction is greater than
$(2\theta+1) . \dfrac{n}{mm}\left\{1-\dfrac{1}{2(2\theta+1)}\dfrac{mm}{nn}\right\}$ or
$\dfrac{n}{mm}(2\theta+1)-\dfrac{1}{2n}$. 
It follows that $P_\theta, P_{\theta+1}, P_{\theta+2}$ and\,so\,on form a strictly
decreasing infinite sequence of quantities provided $2\theta+1$ is at least
$\dfrac{mm}{n}+\dfrac{mm}{2nn}$. Thus it is impossible for all the terms of the sequence
$P, P_1, P_2, P_3$ and\,so\,on to be proportional to integers. Yet this would
necessarily be the case if $m$ and $n$ were integers and
the first two terms of that sequence, $P, P_1$, stood in a rational
ratio to one another. Indeed, we have

\[
\begin{aligned}
P_2 &= nP_1 - mP\\
P_3 &= 3nP_2 - mmP_1\\
P_4 &= 5nP_3 - mmP_2\\
P_5 &= 7nP_4 - mmP_3\\
&\text{and\,so\,on,}
\end{aligned}
\]
\noindent and in general, from the second formula onward,

\[
P_{\theta+2} = (2\theta+1)nP_{\theta+1} - mmP_\theta.
\]

Clearly in these formulas the quantities $P, P_1, P_2$ and\,so\,on can be
replaced by proportional quantities. If the first two of these were integers,
all the subsequent ones would also be integers, which has just been shown to be impossible.

Therefore, when $m, n$ are integers, the value of the fraction
$\dfrac{P_1}{P}$, which is the tangent of the arc
$\dfrac{m}{n}$, cannot be rational.
